\section{Timing of Camera Exposure and Terminal Position}
\label{sec:window552}
This additional study, on which the main paper does not rely, asks whether it matters when a camera perturbation is seen. In $\pi_{0.5}$, moving an equal-length exposure window from an earlier stage to just before clean completion changes neither final position nor success detectably: the primary difference in terminal distance is $-2.925$ mm $[-9.833,\,+1.776]$ and in success $0.000$ points (Table~\ref{tab:supp-window552-estimates}). Exposure for the whole episode does lower success: the difference from the clean run is $-7.971$ points $[-18.841,\,-0.725]$, in a secondary comparison in which the clean run always ran first.

\paragraph{Design: four runs per group, compared within groups.}
The study uses 23 tasks (11 single-support, 12 single-containment), two fixed start states (24 and 25) and three shared policy-noise samples per task and state: 138 matched groups of four runs, 552 episodes. Each group starts with a clean run $C$, which sets the exposure windows. The other three runs apply a $+30^\circ$ camera rotation for the whole episode ($F_{30}$), near clean completion ($R_{30}$), or during the earlier window ($N_{30}$); the attached headlight follows the camera. The subscript is the rotation in degrees. The runs follow a preallocated order. The primary comparison is the task-equal mean of the three paired $R_{30}-N_{30}$ differences at start state 24 on all 69 assigned groups; start state 25 repeats the rule on another 69 groups, and the pooled 138-group analysis is secondary. Pointwise 95\% intervals come from 10,000 shared whole-task bootstrap draws that keep the paired groups within each task.

\paragraph{Timing does not detectably matter; persistent exposure costs success.}
Observed successes for C/$F_{30}$/$R_{30}$/$N_{30}$ are 69/61/68/68 of 69 at state 24 and 68/65/65/66 of 69 at state 25. The terminal-distance intervals span zero in the primary, replication and pooled analyses (Table~\ref{tab:supp-window552-estimates}). $F_{30}-C$ is a different contrast: C always runs first and sets the timing, so this ordered comparison of persistent exposure differs from the randomized comparison of $R_{30}$ and $N_{30}$.

\paragraph{The clean run sets both windows.}
Both windows span four complete policy queries, 20 executed controls. Let $T$ count the clean run's executed controls and $e=5\lfloor T/5\rfloor$. The completion window covers $[e-20,e)$. For the earlier window, let $g$ be the zero-based index of the first transition of the gripper command from nonpositive to positive. Among windows of 20 executed controls aligned to five-step boundaries between $g+5$ and $e-25$, we choose the one whose midpoint is closest to the midpoint of that range, breaking ties toward the earlier window. If none exists and $g\ge45$, the fallback is $[5,25)$, provided it does not overlap the completion window. Both windows require a successful clean run and four complete reference queries; otherwise both assigned runs keep the unchanged camera. Query indices are fixed before the perturbed runs and are not retimed as those runs evolve.

\paragraph{Endpoints: terminal position and task success.}
The position endpoint is measured after each run's last executed control, at its own first success, done signal or fixed horizon: the world-XY distance from the named target's root-body origin to the \emph{current} origin of the named receiver body, or to the declared center of the receiver site, evaluated whether or not the task succeeded. A positive $R_{30}-N_{30}$ means R ends farther from that center; for task success, a positive $\mathrm{SR}(R_{30})-\mathrm{SR}(N_{30})$ means \emph{more} R successes. A run that stops early stays in its assigned group even if this shortens its exposure.

\begin{table*}[t]
\caption{\textbf{Moving the exposure window does not detectably change terminal distance or success; persistent exposure lowers success.} All planned groups (69 per state, 138 pooled), no missing run. Task-equal means of paired within-group differences, left minus right; distances in millimetres; success in percentage points; 95\% intervals from 10,000 whole-task draws. State 24 $R_{30}-N_{30}$ is primary.}

\centering\small
\setlength{\tabcolsep}{5pt}
\begin{tabular}{p{0.35\textwidth}p{0.30\textwidth}p{0.27\textwidth}}
\toprule
Population and contrast & Terminal XY distance difference, mm [95\% CI] & Task success difference, pp [95\% CI]\\
\midrule
State 24 $R_{30}-N_{30}$, primary; 23 tasks, 69/69 pairs & $-2.925\;[-9.833,\,+1.776]$ & $0.000\;[-4.348,\,+4.348]$\\
State 25 $R_{30}-N_{30}$, fixed-state replication; 69/69 pairs & $-1.601\;[-6.788,\,+2.317]$ & $-1.449\;[-8.696,\,+4.348]$\\
Both states $R_{30}-N_{30}$, secondary; 138/138 pairs & $-2.263\;[-8.068,\,+1.640]$ & $-0.725\;[-6.522,\,+4.348]$\\
Both states $F_{30}-C$, secondary; 138/138 pairs & $+1.488\;[-2.914,\,+5.533]$ & $-7.971\;[-18.841,\,-0.725]$\\
\bottomrule
\end{tabular}

\label{tab:supp-window552-estimates}
\end{table*}

\paragraph{All groups stay in the analysis, including no-op and shortened exposures.}
Failures and time limits are observed outcomes; there are no technical failures. The 552 unique episodes cover every planned group, with terminal position and success usable in all 552 episodes, so both intention-to-treat (ITT) contrasts are identified. Each start state has 58 groups in which both windows are defined and 11 in which neither is, none of unknown status. In those 11 groups both R and N apply the declared zero-angle no-op and stay in the denominator; their observed R--N contrasts are zero, and the analysis keeps their actual outcomes rather than imposing equality. Of the 58 eligible groups, N receives all four scheduled queries and 20 executed controls in 58/58 at each state, R in 52/58 at state 24 and 51/58 at state 25; the remaining six and seven R runs have shorter exposure and stay in the ITT analysis. The estimates and exposure records cover all 138 assigned groups; the exploratory subset of 38 groups and the paired executions of the main paper are separate populations; neither replaces this analysis.

\paragraph{Implementation note.}
The bootstrap uses a fixed seed (20261008).
